\documentclass[10pt, letterpaper]{article}

% Packages:
\usepackage[
    ignoreheadfoot, % set margins without considering header and footer
    top=2 cm, % seperation between body and page edge from the top
    bottom=2 cm, % seperation between body and page edge from the bottom
    left=2 cm, % seperation between body and page edge from the left
    right=2 cm, % seperation between body and page edge from the right
    footskip=1.0 cm, % seperation between body and footer
    % showframe % for debugging 
]{geometry} % for adjusting page geometry
\usepackage{titlesec} % for customizing section titles
\usepackage{tabularx} % for making tables with fixed width columns
\usepackage{array} % tabularx requires this
\usepackage[dvipsnames]{xcolor} % for coloring text
\definecolor{primaryColor}{RGB}{0, 0, 0} % define primary color
\usepackage{enumitem} % for customizing lists
\usepackage{fontawesome5} % for using icons
\usepackage{amsmath} % for math
\usepackage[
    pdftitle={Flavio Romano's CV},
    pdfauthor={Flavio Romano},
    colorlinks=true,
    urlcolor=primaryColor
]{hyperref} % for links, metadata and bookmarks
\usepackage[pscoord]{eso-pic} % for floating text on the page
\usepackage{calc} % for calculating lengths
\usepackage{bookmark} % for bookmarks
\usepackage{lastpage} % for getting the total number of pages
\usepackage{changepage} % for one column entries (adjustwidth environment)
\usepackage{paracol} % for two and three column entries
\usepackage{ifthen} % for conditional statements
\usepackage{needspace} % for avoiding page brake right after the section title
\usepackage{iftex} % check if engine is pdflatex, xetex or luatex

% Ensure that generate pdf is machine readable/ATS parsable:
\ifPDFTeX
    \input{glyphtounicode}
    \pdfgentounicode=1
    \usepackage[T1]{fontenc}
    \usepackage[utf8]{inputenc}
    \usepackage{lmodern}
\fi

\usepackage{charter}

% Some settings:
\raggedright
\AtBeginEnvironment{adjustwidth}{\partopsep0pt} % remove space before adjustwidth environment
\pagestyle{empty} % no header or footer
\setcounter{secnumdepth}{0} % no section numbering
\setlength{\parindent}{0pt} % no indentation
\setlength{\topskip}{0pt} % no top skip
\setlength{\columnsep}{0.15cm} % set column seperation
\pagenumbering{gobble} % no page numbering

\titleformat{\section}{\needspace{4\baselineskip}\bfseries\large}{}{0pt}{}[\vspace{1pt}\titlerule]

\titlespacing{\section}{
    % left space:
    -1pt
}{
    % top space:
    0.15 cm
}{
    % bottom space:
    0.1 cm
} % section title spacing

\renewcommand\labelitemi{$\vcenter{\hbox{\small$\bullet$}}$} % custom bullet points
\newenvironment{highlights}{
    \begin{itemize}[
        topsep=0.10 cm,
        parsep=0.10 cm,
        partopsep=0pt,
        itemsep=0pt,
        leftmargin=0 cm + 10pt
    ]
}{
    \end{itemize}
} % new environment for highlights


\newenvironment{highlightsforbulletentries}{
    \begin{itemize}[
        topsep=0.10 cm,
        parsep=0.10 cm,
        partopsep=0pt,
        itemsep=0pt,
        leftmargin=10pt
    ]
}{
    \end{itemize}
} % new environment for highlights for bullet entries

\newenvironment{onecolentry}{
    \begin{adjustwidth}{
        0 cm + 0.00001 cm
    }{
        0 cm + 0.00001 cm
    }
}{
    \end{adjustwidth}
} % new environment for one column entries

\newenvironment{twocolentry}[2][]{
    \onecolentry
    \def\secondColumn{#2}
    \setcolumnwidth{\fill, 4.5 cm}
    \begin{paracol}{2}
}{
    \switchcolumn \raggedleft \secondColumn
    \end{paracol}
    \endonecolentry
} % new environment for two column entries

\newenvironment{threecolentry}[3][]{
    \onecolentry
    \def\thirdColumn{#3}
    \setcolumnwidth{, \fill, 4.5 cm}
    \begin{paracol}{3}
    {\raggedright #2} \switchcolumn
}{
    \switchcolumn \raggedleft \thirdColumn
    \end{paracol}
    \endonecolentry
} % new environment for three column entries

\newenvironment{header}{
    \setlength{\topsep}{0pt}\par\kern\topsep\centering\linespread{1.5}
}{
    \par\kern\topsep
} % new environment for the header

\newcommand{\placelastupdatedtext}{% \placetextbox{<horizontal pos>}{<vertical pos>}{<stuff>}
  \AddToShipoutPictureFG*{% Add <stuff> to current page foreground
    \put(
        \LenToUnit{\paperwidth-2 cm-0 cm+0.05cm},
        \LenToUnit{\paperheight-1.0 cm}
    ){\vtop{{\null}\makebox[0pt][c]{
        \small\color{gray}\textit{Last updated in August 2026}\hspace{\widthof{Last updated in August 2026}}
    }}}%
  }%
}%

% save the original href command in a new command:
\let\hrefWithoutArrow\href

% new command for external links:


\begin{document}
\newcommand{\AND}{\unskip
    \cleaders\copy\ANDbox\hskip\wd\ANDbox
    \ignorespaces
}
\newsavebox\ANDbox
\sbox\ANDbox{$|$}

\begin{header}
    \fontsize{25 pt}{25 pt}\selectfont Flavio Romano

    \vspace{5 pt}

    \normalsize
    \mbox{Italy}%
    \kern 5.0 pt%
    \AND%
    \kern 5.0 pt%
    \mbox{\hrefWithoutArrow{mailto:flavromano@proton.me}{flavromano@proton.me}}%
    \kern 5.0 pt%
    \AND%
    \kern 5.0 pt%
    \mbox{\hrefWithoutArrow{https://www.linkedin.com/in/flavromano/}{linkedin.com/in/flavromano}}%
    \kern 5.0 pt%
    \AND%
    \kern 5.0 pt%
    \mbox{\hrefWithoutArrow{https://github.com/FlavRomano}{github.com/FlavRomano}}%
\end{header}

\vspace{5 pt - 0.3 cm}

\vspace{.4 cm}

Software Engineer and Technical Lead specializing in backend systems, identity \& access management, and cloud-native microservices architectures. Experienced in designing distributed systems, leading engineering teams, and delivering enterprise-grade software in telecommunications and sustainability domains. Currently pursuing an MSc in Artificial Intelligence, adding advanced AI/ML expertise to a strong backend engineering foundation.
\section{Experience}

\begin{twocolentry}{
        March 2024 – Present
    }
    \textbf{Software Engineer / Technical Lead}, NTT DATA -- Pisa, IT
\end{twocolentry}

\vspace{0.10 cm}
\begin{onecolentry}
    \begin{highlights}

        \item Contributed to the design and implementation of OAuth 2.0, OpenID Connect, and SAML authentication solutions for a leading Italian telecommunications provider, enhancing a customized Keycloak deployment and enabling adoption of the Phantom Token architecture.

        \item Developed Spring Boot microservices supporting authorization workflows, including token introspection, redirect management, and JWT/JWE processing.

        \item Served as Technical Lead for the development of an AI-powered platform automating due diligence documentation for European sustainability regulations, coordinating a team of 4 engineers and leading technical decisions.

        \item Currently leading a team of 5 engineers in the development of a sustainability platform for estimating CO$_2$ emissions generated by software systems, hardware infrastructure, and AI workloads, using Spring WebFlux for the reactive CO$_2$ estimation service, PostgreSQL for customer data, and Grafana/ELK for monitoring and log analysis.
    \end{highlights}
\end{onecolentry}

\vspace{0.1 cm}

\begin{twocolentry}{
        Dec 2023 – March 2024
    }
    \textbf{Software Engineer Intern}, NTT DATA -- Pisa, IT
\end{twocolentry}

\vspace{0.10 cm}
\begin{onecolentry}
    \begin{highlights}
        \item Developed a Java microservice for the GreenIT ecosystem to collect hardware metrics and estimate Software Carbon Intensity (SCI) and CO$_2$ emissions on Google Cloud Platform infrastructures.

        \item Designed a configurable monitoring connector that later became the foundation of my Bachelor's thesis.
    \end{highlights}
\end{onecolentry}


\section{Technical Skills}
\begin{onecolentry}
    \textbf{Languages:} Java, Python, Go, TypeScript, SQL, Bash, C \\
    \textbf{Backend \& Cloud:} Spring Boot, Spring WebFlux (Reactive), FastAPI, PostgreSQL, Kubernetes, AWS, GCP, RabbitMQ, Redis, Terraform, Grafana, ELK Stack \\
    \textbf{Security:} OAuth 2.0, OpenID Connect, SAML, Keycloak \\
    \textbf{ML:} PyTorch, Scikit-learn, Keras, Pandas, Optuna, RAG, LLM Evaluation
\end{onecolentry}


\section{Projects}

\vspace{0.1 cm}
\begin{twocolentry}{
        Apr 2026 – May 2026
    }
    \textbf{Encoding Classification using LSTM} (\href{https://github.com/FlavRomano/encoding-classifier-lstm}{github})
\end{twocolentry}

\begin{highlights}
    \item Developed a PyTorch-based sequence classification model to automatically identify the encoding algorithm (Caesar Cipher, ROT13, MD5) applied to a text string.

    \item Built and trained a 230k-parameter LSTM on 72,000+ samples, achieving 97.99\% test accuracy, with automated training/inference workflows for real-time prediction.
\end{highlights}

\begin{twocolentry}{
        Aug 2023 – Dec 2023
    }
    \textbf{Coinquilining} (\href{https://github.com/FlavRomano/coinquilining}{github})
\end{twocolentry}

\begin{highlights}
    \item Built from scratch as final project for the "Web App Development" course (top grade, 30/30): an installable, offline-capable PWA for shared households, covering fridge/pantry inventory, shopping lists, cleaning rotations, and expense splitting/balances.

    \item Implemented expiration-tracking notifications, ingredient-based recipe search via Puppeteer scraping, and a fair cleaning-rotation algorithm (time shift problem) rendered as a drag-and-drop calendar. Tech Stack: SvelteKit, TypeScript, Supabase (Postgres/Auth), Tailwind CSS, Puppeteer
\end{highlights}


\section{Education}

\begin{twocolentry}{
        Sept 2025 – Present
    }
    \textbf{University of Pisa}, Master of Science (MSc) in Artificial Intelligence \end{twocolentry}

\vspace{0.10 cm}
\begin{twocolentry}{
        Sept 2020 – May 2024
    }
    \textbf{University of Pisa}, Bachelor of Computer Science\end{twocolentry}
\end{document}